\section{Reconciliation Certificates}\label{sec:cert}

Reconciliation is the verifier (\adr{004}). It checks equivalence on data
rather than on syntax, so it applies to changes that no SQL prover accepts, and
it states exactly which keys it checked.

\subsection{The verifier}\label{sec:verifier}

Reconciliation compares a left output $L$ (the baseline) with a right output
$R$ (the candidate) by key. Its parameters are the key columns, an optional
latest-by column, ignored columns, and group and decimal columns for totals.
For each key $k$ in a scope $Q$ it takes the latest-version rows of $k$ on each
side, $L_k$ and $R_k$ (the rows at $k$'s greatest latest-by value; null
timestamps never win, \adr{014}), and compares them by a row hash $h$: xxh3-128
from the xxHash family~\cite{xxhash} over the canonical encodings of the
non-ignored columns in schema order, each preceded by its column name's length
and bytes, with an explicit null marker (\doc{format.md}, \adr{016}). It also
sums decimal columns per group exactly, at the dataset's declared scale. The
Certificate records
\begin{align*}
  \mathit{missing\_right} &= |\{ k \in Q : L_k \neq \emptyset,\ R_k = \emptyset \}|,\\
  \mathit{missing\_left}  &= |\{ k \in Q : L_k = \emptyset,\ R_k \neq \emptyset \}|,\\
  \mathit{mismatched}     &= |\{ k \in Q : L_k, R_k \neq \emptyset,\
                              h(L_k) \neq h(R_k) \}|,
\end{align*}
the ignored columns, and every group whose totals differ, with both totals at
the declared scale. It passes if and only if the three counts are zero and no
totals differ (\doc{agent-loop.md}). The index makes this cheap (\adr{004}):
the scope of a workload certificate is the set of keys the recorded runs
touched, and the index locates them to row groups, so running the job both ways
and reconciling reads the row groups that hold touched keys rather than the
dataset.

\subsection{Coverage}\label{sec:coverage}

Every Certificate states its coverage: \emph{full} (every key),
\emph{workload} (the keys of a named recorded workload, with their count), or
\emph{sampled} (a number of keys and the sampled fraction).
\emph{Guarantee (design property):} a passing Certificate with scope $Q$
asserts that for every $k \in Q$ the latest-version rows of $k$ agree on every
non-ignored column, up to collisions of a 128-bit non-cryptographic hash, and
that group totals over the compared rows agree exactly. It asserts nothing
about keys outside $Q$; in the design's words, a Certificate is a statement
about the keys it covers and nothing more (\doc{agent-loop.md}).

\subsection{Why provers do not apply}\label{sec:provers}

SQL equivalence provers and deciders establish equivalence for all database
instances, or for all instances up to a bound, within a supported SQL
fragment~\cite{cosette,qed,sqlsolver,verieql}. Bellwether's changes fall outside
that setting in three ways (\adr{004}). Pipeline steps are written as Spark
plans, KQL, Scala, or configuration, which these tools do not accept. Many
changes are not query rewrites at all: a key spec, a bloom false-positive rate,
a lookback bound, a layout. And the rewrites with the largest effect often cross
step boundaries and change intermediate results
(Section~\ref{sec:nonlocal}). Reconciliation applies to all three because it
compares outputs; the price is that a Certificate covers the keys it compared,
not all inputs. The two kinds of evidence compose: where a rule catalog identity
applies, such as a semi-join on $K$ commuting with a window partitioned by $K$,
the gate may use the identity for that edit and reconciliation for the rest
(\doc{evolve.md} \S1, \S4.2). In the current contract a static identity has no
Certificate path and so cannot reach a Receipt; the planned \code{Verifier}
trait adds one (\doc{rule-catalog.md}; \doc{evolve.md} \S10).

\subsection{Whole-pipeline equivalence and non-local rewrites}\label{sec:nonlocal}

A pipeline is a DAG of steps with lineage. The rewrites that matter most cross
step boundaries: push a filter from a late step to an early one, materialize an
intermediate that two steps share, or remove a deduplication that a later step
repeats (\doc{evolve.md} \S5.3). Each changes some step's local output, so a
check that requires equivalence step by step must reject it. Bellwether gates at
the pipeline output instead: for a candidate $a$ and the baseline $a_0$,
$\Eq(a, a_0)$ holds when reconciliation of the two pipeline outputs over the
touched keys passes, and intermediate outputs may differ (\adr{023}).

\emph{Illustration.} Step two of a pipeline reads the full history of a
versioned table and keeps the newest version of each key; step five semi-joins
that result with the run's trigger keys. The candidate moves the semi-join into
step two and answers it with a temporal \emph{latest} lookup. Step two's output
now holds only the trigger keys, so a per-step check fails, yet step five's
output is unchanged for every trigger key, which is what the output certificate
checks. The candidate composes two catalog patterns,
\code{lookback\us unbounded} and \code{unprefiltered\us window}
(\doc{rule-catalog.md}), across a step boundary;
Figure~\ref{fig:schematic-pipeline} draws it.

\subsection{Limits}\label{sec:certlimits}

A non-local rewrite is certified for the keys the workload touched, not for all
inputs, so a rewrite that is wrong only on untouched keys would pass. How to
communicate this to a user who later runs an untouched key is an open question
in the design (\doc{evolve.md} \S13), and it is why a rule must be re-validated
on held-out runs before promotion (Section~\ref{sec:rulebook}). A sampled
certificate speaks only for its sample, and ignored columns are not compared,
which is why the Certificate lists them. The row hash is not cryptographic, so a
Certificate does not defend against an adversary who can craft colliding rows;
an attacker with write access to the store is already outside the threat model.
Data can also drift: each decision records the fanout evidence it depended on,
and a receipt whose evidence has drifted is flagged stale (\doc{evolve.md}
\S5.4).
